\chapter{Verificación del simulador}\label{ap:verificacion}

Antes de generar los resultados se verificó que el simulador no lineal fuera coherente con el modelo de diseño. La verificación detectó dos errores en una versión previa de la función de la planta: la fuerza electromagnética se dividía dos veces entre la masa, y la fricción viscosa se calculaba como $c_1|v|$, por lo que no se oponía al movimiento cuando la velocidad era negativa. Con esos errores, el voltaje de equilibrio en $x_1=-4$~cm resultaba de $0{,}49$~V en lugar de $3{,}48$~V, y el voltaje de control quedaba siempre muy por debajo del límite del actuador.

Tras la corrección se comprobaron tres propiedades.

\begin{itemize}
  \item El voltaje que anula la fuerza neta en $x_1=-4$~cm es $3{,}4826$~V, igual al voltaje de equilibrio del modelo linealizado.
  \item El jacobiano de la función de la planta, calculado numéricamente en $x_1=-4$~cm, coincide con la matriz $A$ del modelo linealizado, con valores propios $14{,}9427$ y $-23{,}5057$.
  \item La fricción viscosa tiene signo opuesto a la velocidad: la aceleración que produce es $-8{,}563$~cm/s$^2$ con $v=+1$~cm/s y $+8{,}563$~cm/s$^2$ con $v=-1$~cm/s.
\end{itemize}

La corrección cambió de forma importante el comportamiento del lazo. Con el modelo corregido, el voltaje necesario para operar en la región de $-4$ a $-5$~cm se acerca al límite del actuador, lo que motivó el estudio del intervalo de operación de la sección~\ref{sec:intervalo} y la elección del intervalo de $-3$ a $-2$~cm para las pruebas.
